\documentclass[11pt,a4paper]{article}
\usepackage[margin=2.0cm]{geometry}
\usepackage[T1]{fontenc}
\usepackage{lmodern}
\usepackage{microtype}
\usepackage{parskip}
\usepackage{enumitem}
\usepackage{fancyhdr}
\usepackage[hidelinks]{hyperref}

\pagestyle{fancy}
\fancyhf{}
\fancyhead[L]{\small NEUCOM-D-26-24287R1 --- Cover letter (second revision)}
\fancyhead[R]{\small\thepage}
\renewcommand{\headrulewidth}{0.4pt}

\setlist[itemize]{topsep=2pt,itemsep=1.5pt,parsep=0pt,leftmargin=1.3em}

\begin{document}

\begin{flushright}
\small
Jiang Tao \quad jiangtao.ml@foxmail.com\\
School of Intelligent Information Engineering\\
Heilongjiang University of Finance and Economics, Harbin, China
\end{flushright}

\vspace{-6pt}
\noindent\small \today

\vspace{7pt}
\noindent The Editor-in-Chief\\
\noindent\emph{Neurocomputing}

\vspace{7pt}
\noindent\textbf{Re: NEUCOM-D-26-24287R1 --- second revision of ``FOLoRA: Fisher-Orthogonalized
Low-Rank Adaptation of One Shared Adapter for Continual Learning''}

\vspace{6pt}
Dear Editor,

Two reviewers asked for changes in this round, and both asked for something that alters what
the paper claims. We have made both changes rather than argue around them, and in each case
the consequence is that the paper claims less than the version you reviewed.

\medskip
\noindent\textbf{The main table now runs O-LoRA at the objective its own paper specifies.}
Reviewer~1 is right that our principal comparison had been against an initialization-only
variant while every other baseline ran at its full objective. O-LoRA now appears at its
training-time orthogonality constraint ($\lambda_1{=}1.0$), and the initialization-only
variant is demoted to an ablation. This weakens our result: the margin over O-LoRA falls from
$+12.92$ to $+5.60$ points on CIFAR-100 and from $+13.20$ to $+5.74$ on ImageNet-R, and on
ImageNet-R the average-accuracy difference no longer clears Holm correction
($p_{\mathrm{Holm}}{=}0.052$). Because the confirmatory family changed, every corrected
$p$-value was recomputed and the results text, abstract, highlights and conclusion were
updated together.

\medskip
\noindent\textbf{The paper no longer attributes the margin to the Fisher weighting.} Reviewer~8
draws the conclusion our own revision-1 experiments support: since a tuned isotropic
displacement penalty reaches the same accuracy, what separates FOLoRA from the
per-task-module family is that the adapter is \emph{shared}, not that the penalty is
Fisher-weighted. The title, abstract and highlights now say this; the isotropic penalty is a
row of the main comparison rather than a line in the limitations; and Section~5.6 and the
conclusion state that, at this budget, we could not establish any accuracy or forgetting
advantage for the Fisher eigendecomposition --- which also costs $182.8$\,MiB of retained state
where the isotropic penalty costs one scalar per layer.

\medskip
\noindent\textbf{The other items.} Section~4 is $441$ words shorter than the version you
reviewed ($-13.9\%$) with the survey of prior bounds moved to Appendix~C; the abstract now
states the three scope limitations; the learned-head anomaly is explained with a per-tier
recency diagnostic rather than corrected; ``replay-free'' is qualified as ``replay-free during
training''; the TOST margin is justified independently of the observed differences; the term
``pre-registered'' is kept only where a dated record exists, and that record is quoted verbatim
in the response; a validation-split re-selection of $\lambda$ is reported; and the five
unrelated metaheuristic references are removed.

\medskip
\noindent\textbf{Files.} The revision is accompanied by the response to reviewers, the revised
manuscript (marked and unmarked), the \LaTeX{} source, the figure files, updated highlights and
abstract, the per-seed results behind every number including the new validation experiment, and
a one-page evidence log. Declarations are unchanged.

\vspace{1pt}
\noindent Sincerely,\\[-2pt]
Jiang Tao\\
Corresponding author

\end{document}
